\documentclass[11pt]{article}
\pdfobjcompresslevel=0 % force classic uncompressed xref table (avoids MiKTeX broken-xref-stream bug)
\usepackage[margin=2.3cm]{geometry}
\usepackage[T1]{fontenc}\usepackage{lmodern}
\usepackage{amsmath,amssymb,mathtools}
\usepackage{tcolorbox}\tcbuselibrary{skins,breakable}
\usepackage{booktabs}\usepackage{array}\usepackage{enumitem}
\usepackage{fancyhdr}\setlength{\headheight}{14.5pt}
\usepackage{parskip}\usepackage{xcolor}\usepackage{hyperref}
\usepackage{tikz}\usetikzlibrary{arrows.meta,positioning}

\definecolor{dtfblue}{HTML}{1a1a2e}\definecolor{dtflight}{HTML}{0f3460}
\definecolor{boxblue}{HTML}{e8f4f8}\definecolor{knownblue}{HTML}{1a3a6e}
\definecolor{cbKnownBg}{HTML}{E9E9FA}\definecolor{cbKnownFr}{HTML}{8E97C9}
\definecolor{cbNodeBg}{HTML}{F4F5F8}\definecolor{cbNodeFr}{HTML}{CDD2DB}
\definecolor{cbAssumeBg}{HTML}{FBEBEB}\definecolor{cbAssumeFr}{HTML}{C0392B}
\definecolor{cbResultBg}{HTML}{FBF4D2}\definecolor{cbResultFr}{HTML}{C8A21E}
\definecolor{cbTagBg}{HTML}{E2F1EF}\definecolor{cbTagFr}{HTML}{4FA295}
\definecolor{cbArrow}{HTML}{8A93A6}\definecolor{cbPyBg}{HTML}{ECEFF4}

\tcbset{
  keybox/.style={enhanced,breakable,colback=boxblue,colframe=dtfblue,boxrule=0.8pt,
    arc=3pt,left=10pt,right=10pt,top=8pt,bottom=8pt},
  cbflow/.style={enhanced,arc=3.5pt,boxrule=0.7pt,left=12pt,right=12pt,top=10pt,bottom=10pt,
    before skip=4pt,after skip=4pt},
}
\newcommand{\chainarrow}{\vspace{1pt}\centerline{\tikz{\draw[-{Latex[length=2.6mm,width=2.2mm]},
  line width=1pt,cbArrow](0,0.28)--(0,0);}}\vspace{1pt}}
\newcommand{\known}[2]{\begin{tcolorbox}[cbflow,colback=cbKnownBg,colframe=cbKnownFr]
  \centering\(\displaystyle #1\)\\[3pt]{\color{knownblue}\small[\,\textit{known: #2}\,]}\par
  \nobreak\vspace{1pt}\hfill{\tiny\color{gray}Known}\end{tcolorbox}}
\newcommand{\assume}[1]{\chainarrow\begin{tcolorbox}[cbflow,colback=cbAssumeBg,colframe=cbAssumeFr]
  {\color{cbAssumeFr}\bfseries[DTF input:]\ }{\color{cbAssumeFr}\small #1}\par
  \nobreak\vspace{1pt}\hfill{\tiny\color{gray}DTF input}\end{tcolorbox}}
\newcommand{\nodetext}[1]{\chainarrow\begin{tcolorbox}[cbflow,colback=cbNodeBg,colframe=cbNodeFr]
  \small #1\par\nobreak\vspace{1pt}\hfill{\tiny\color{gray}Derivation}\end{tcolorbox}}
\newcommand{\dtfresult}[2]{\chainarrow
  \begin{tcolorbox}[cbflow,colback=cbResultBg,colframe=cbResultFr,boxrule=1.3pt]
  \centering\(\displaystyle #1\)\par\nobreak\vspace{1pt}\hfill{\tiny\color{gray}DTF result}\end{tcolorbox}\chainarrow
  \begin{tcolorbox}[cbflow,colback=cbTagBg,colframe=cbTagFr]
  \centering{\itshape\color{dtflight}#2}\par\nobreak\vspace{1pt}\hfill{\tiny\color{gray}Takeaway}\end{tcolorbox}}
\newcommand{\pychip}[2]{{\footnotesize\colorbox{cbPyBg}{\ttfamily\,#1\,}\,{\color{gray}$\to$ #2}}}
\newcommand{\resultpy}[3]{\chainarrow
  \begin{tcolorbox}[cbflow,colback=cbResultBg,colframe=cbResultFr,boxrule=1.3pt]
  \centering\(\displaystyle #1\)\par\smallskip\centering\pychip{#2}{#3}\par
  \nobreak\vspace{1pt}\hfill{\tiny\color{gray}DTF result --- verified}\end{tcolorbox}}
\newcommand{\oneobject}[2]{\begin{tcolorbox}[enhanced,colback=white,colframe=black,boxrule=1.3pt,arc=3pt,
  left=10pt,right=10pt,top=7pt,bottom=7pt,before skip=4pt,after skip=4pt]
  {\bfseries One object, seen here as: #1}\\[2pt]{\small #2}\end{tcolorbox}}
\newenvironment{chainbox}{\par\addvspace{4pt}\noindent\ignorespaces}{\par\addvspace{4pt}}
\newcommand{\chainhead}[2]{\subsection*{Chain #1\quad\normalfont\itshape #2}}
\newcommand{\setup}[1]{\noindent\emph{#1}\par\smallskip}
\newcommand{\tr}{\operatorname{tr}}

\hypersetup{colorlinks=true,linkcolor=knownblue,citecolor=knownblue,urlcolor=knownblue,
  pdftitle={The Distorted Time Framework, Volume III: The Definite Clock},
  pdfauthor={William de Dufour}}
\pagestyle{fancy}\fancyhf{}
\fancyhead[L]{\small\itshape Distorted Time Framework --- Volume III}
\fancyhead[R]{\small\itshape The Definite Clock}
\fancyfoot[C]{\small\thepage}\renewcommand{\headrulewidth}{0.4pt}

\begin{document}

\begin{center}
{\large\bfseries The Distorted Time Framework\\[3pt]
\normalsize Volume III --- \itshape The Definite Clock}\\[6pt]
{\small The Born rule, measurement, and entanglement from a clock that stays definite}\\[10pt]
{\normalsize William de Dufour}\\[2pt]
{\small\itshape Independent Researcher, Austin, TX}\\[2pt]
{\small\texttt{billd.dtf@gmail.com}}\\[2pt]
{\small ORCID: \href{https://orcid.org/0009-0008-3470-1231}{0009-0008-3470-1231}}\\[2pt]
{\small\url{https://github.com/billd-dtf/Distorted-Time-Framework}}\\[3pt]
{\small Version 8}\\[3pt]
{\small\href{https://doi.org/10.5281/zenodo.21659853}{\texttt{DOI: 10.5281/zenodo.21659853}}}
\end{center}

\begin{abstract}
Quantum mechanics postulates the Born rule and has never derived it; the closest route --- Nelson's
stochastic mechanics --- has stalled on Wallstrom's single-valuedness obstruction and on the unexplained
origin of its noise. Taking time's local rate as primitive (Volume~I), this volume reads a quantum's phase
as the clock's reading and \emph{reduces} the Born rule to the relaxation equilibrium of the conserved
location density under the native guiding flow. Wallstrom is reduced to a single \emph{declared} posit ---
that Row~0's beable is a single-valued amplitude assignment --- for the simply-connected sector, the nodal
case left open. Nelson's missing noise is named: the internal flexing (a zitterbewegung) of bound energy at
the Compton scale, fixing the diffusion coefficient to $\hbar/2m$; its conservative (time-symmetric)
character --- hence Schr\"odinger rather than a heat equation --- follows from the flexing being a
\emph{bound} mode, which cannot dissipate, on a timeless substrate. Entanglement is a shared present on the
proper now, with $\mathrm{CHSH}=2\sqrt2$ inherited, not derived. Two honest limits are kept in view: the
guiding wave lives on configuration space ($O(N^2)$ content against a slice's $O(N)$), so the local-beable
reading holds only for one particle; and the amplitude sector is non-relativistic, its relativistic
completion (toward Dirac) an open program. The claim throughout is reduction and identification, not a
derivation of quantum mechanics.
\end{abstract}

\smallskip
\noindent{\small\textbf{Keywords:} Born rule; Nelson stochastic mechanics; Wallstrom obstruction;
pilot-wave theory; zitterbewegung; entanglement; measurement.}

\bigskip

\begin{tcolorbox}[keybox]
\textbf{The three commitments (from the Core, Volume~I).}\enspace This volume builds on three, argued there
and used here without re-argument (full recap: Appendix~A):
\begin{itemize}[leftmargin=1.4em,itemsep=1pt,topsep=3pt]
  \item \textbf{1.\ Time is primitive} --- its local rate $u$ is the one physical thing, and a quantum's
    phase is its reading, $\Phi=-\omega_C\!\int u\,dt$ ($\omega_C=mc^2/\hbar$): Row~0 (timeless,
    $\hat H\Psi=0$) $\to$ Row~1 (the definite clock) $\to$ Row~2 (space); ``becoming definite'' is the
    Row-0$\to$1 step.
  \item \textbf{2.\ Space is time's radiation} --- and the moving clock's momentum is
    $\mathbf p=\nabla S=m\mathbf v$: inertia \emph{is} the rate $\omega_C$, the pilot-wave guidance the same
    space-face gradient $\nabla S$ ($\hbar$-coupled, present even at $u\equiv1$); the $G$-coupled well the
    rate digs is gravity, a separate object.
  \item \textbf{3.\ There is a proper now} --- the preferred slice a non-propagating rate fixes, real but
    empirically silent; here it is the shared present that carries entanglement.
\end{itemize}
\noindent\emph{Here} the rate's reading is the quantum phase, and its residual internal flexing is the
quantum noise.
\end{tcolorbox}

\medskip
\noindent Quantum mechanics postulates the Born rule~\cite{Born1926} and has never derived it. The route
that came closest --- Nelson's stochastic mechanics~\cite{Nelson1966} --- has been held up thirty years by
one obstruction (Wallstrom's~\cite{Wallstrom1994}: nothing forces the phase to be single-valued, so the
quantised spectrum is lost). This volume \textbf{reduces Born to a relaxation equilibrium on three
physical premises} --- and \emph{reduces} the Wallstrom obstruction to a single ontological posit rather
than an imposed quantisation rule: take the fundamental object to be Row~0's amplitude assignment
($\hat H\psi_0=0$), a single-valued function on configuration space, and integer circulation then follows
by topology alone, with nothing added by hand. That posit is \emph{declared}, not derived --- and it is a
genuine restriction, since physically multivalued states exist (spinors, anyons) --- so the reduction is
made for the simply-connected sector and the nodal / multiply-connected case stays open. It keeps de~Broglie--Bohm's~\cite{Bohm1952}
commitment that the particle is real and guided. On the wave equation itself the claim is bounded and
stated once: DTF derives the \emph{phase} half of the Schr\"odinger equation --- the space-face of its own
interval --- and \emph{reduces} the amplitude half to a single posited scale, $\hbar$, the clock's jitter;
it does not claim to derive the equation, and the closing box totals this at its true size.

\noindent It does \emph{not} make the guiding object local. Counting settles this rather than leaving it
open: a general two-particle amplitude needs functional content growing like $N^2$ on $N$ cells of space,
while everything DTF puts on a slice --- the clock rate and the extent structure --- grows like $N$. No
bounded amount of internal structure closes that gap, and enlarging it until it does is configuration space
with an index renamed. So the guiding wave lives on configuration space here as in every Bohmian theory;
the local-beable reading holds for a single quantum and does not extend
[\texttt{DTF\_sharednow\_counting.py}].

\noindent \textbf{What is distinctive survives that, and is the volume's actual claim.} Relativistic
Bohmian mechanics needs a preferred foliation and must postulate one. DTF's is not postulated: it falls
out of $\pi_u\equiv0$, and it is the same slice on which gravity already solves its field --- introduced
in Volume~I for reasons that have nothing to do with Bell. The economy is in where the foliation comes
from, not in making the wave local.

\clearpage
\chainhead{1}{The wave is a phase field, not a vanishing}
\setup{The particle stays definite and stays somewhere. What ``spreads'' is not the particle but the phase field that guides it --- a gradient in space, not a wake in the rate of time.}
\begin{chainbox}
\known{\psi=R\,e^{iS/\hbar},\qquad \Phi=-\omega_C\tau\ \text{(the clock reading)}}{polar form; phase as elapsed proper time}
\known{S=\hbar\Phi=-mc^2\!\int d\tau:\quad \mathbf p=\nabla S=\gamma m\mathbf v\ (\text{space-face}),\quad E=-\partial_t S=\gamma mc^2\ (\text{time-face})}{Hamilton--Jacobi: the one phase has two faces ($\int u\,dt$ at rest)}
\nodetext{a clock carries two independent things: a \emph{rate} (its mass $\omega_C$, which moves it
and bends time around it) and a \emph{state} (a spin). Entanglement touches only the state.}
\resultpy{\text{singlet: spin purity } 0.5,\quad \text{each rest energy sharp, spread } 0.0000}
{DTF\_two\_clocks\_resolved.py}{rate untouched}
\resultpy{u\equiv1\ \Rightarrow\ \nabla S=\gamma m\mathbf v\neq0\ \ (\text{guidance is the space face})}
{DTF\_guidance\_is\_spaceface.py}{no $G$, no $u$-wake}
\dtfresult{\text{the ``wave''}=\text{the phase gradient }\nabla S,\qquad \nabla S\neq0\ \text{even at }u\equiv1}
{The particle never de-materialises. What guides it is the \emph{space}-face phase gradient $\nabla S$ ---
which survives even with the clock rate perfectly uniform ($u\equiv1$), so it is no wake in $u$. The
\emph{rate} itself (the mass $\omega_C$) is the separate thing entanglement leaves razor-sharp.}
\end{chainbox}
\oneobject{inertia \& the pilot wave}{Both read the one phase $\Phi=-\omega_C\tau$, and \emph{neither
carries $G$}. Inertia \emph{is} the rate $\omega_C$ --- the mass $m=\hbar\omega_C/c^2$, linear in $m$,
definitional: the clock's ticking, not a mechanism. Guidance is the \emph{space}-face gradient $\nabla S$
--- $\hbar$-coupled, and it \emph{survives at $u\equiv1$} (a flat clock field). The two are tied because
the momentum \emph{is} that gradient, $\mathbf p=\nabla S=m\mathbf v$: one clock, read as its rate and as
the gradient it lays down. What is \emph{neither} --- the $G$-coupled well the rate digs in $u$,
$\sim(m/m_P)^2$ --- is gravity (Volume~II). Met again through Volumes~I and~II.}

\chainhead{2}{Amplitude is the same clock, read on the space axis}
\setup{The phase read along time is a period; along space it is an extent. One constant ties them.}
\begin{chainbox}
\known{E\,T_C=h\ \ (\text{phase}),\qquad E\,\lambda_C=hc\ \ (\text{amplitude}),\qquad \lambda_C=c\,T_C}
{Planck/Compton, one on each axis (Volume~I, Chain~5)}
\dtfresult{h\ \text{and}\ hc\ \text{differ only by}\ c}
{The phase-anchor and amplitude-anchor are one constant, not two --- ``space is time's radiation.''
The pilot-wave tradition reads the time face, textbook Born the space face; the clock supplies both.}
\end{chainbox}

\begin{tcolorbox}[keybox]
\textbf{The most legible line in the framework: the double slit.}\enspace With $\Phi=-\omega_C\tau$,
bright fringes sit where the two routes' accumulated \emph{proper times} differ by a whole number of
Compton periods --- constructive interference is $\omega_C\,\Delta\tau=2\pi n$, i.e. $\Delta\tau=nT_C$.
For two paths reaching the same lab time this is \emph{exactly} de~Broglie's $\Delta\ell=nh/p$
[\texttt{DTF\_doubleslit\_compton\_ticks.py}, verified exact --- not small-angle]. Label it honestly:
the relativistic action phase rewritten in clock language, a \emph{consistency check}, not new physics.
Its worth is that for \emph{one} particle configuration space \emph{is} ordinary 3-space, so the
single-particle interferometer is where DTF's local-beable reading holds with \textbf{nothing owed} ---
its best case, not its hardest. (Two particles is Chain~4, where the debt begins.)
\end{tcolorbox}

\clearpage
\chainhead{3}{The Born rule as the flow's equilibrium}
\setup{$|\psi|^2$ is not a probability postulate --- it is the density of where the definite particle is. So compute where it settles.}
\begin{chainbox}
\known{\rho=|\psi|^2,\qquad \partial_t\rho+\nabla\!\cdot(\rho\mathbf v)=0}{polar form + local conservation}
\assume{the particle flows along its guiding phase gradient, $\mathbf v=\nabla S/m$ --- the same $\nabla S$ that is its momentum, not a new law}
\nodetext{conservation makes $\rho=|\psi|^2$ the \emph{stationary} distribution of the flow
(equivariance); the settling toward it is a coarse-grained $H$-theorem. The remaining ingredient is
single-valuedness --- Wallstrom's obstruction, that the hydrodynamic pair $(\rho,S)$ does not force
integer circulation,
\[ \oint\nabla S\cdot d\ell\ \in\ 2\pi\hbar\,\mathbb Z. \]
Here it is not imposed as a separate quantisation rule
but \emph{reduced to one ontological posit}: take DTF's fundamental object to be Row~0's amplitude
assignment $\psi_0$ ($\hat H\psi_0=0$), a single-valued function on configuration space, rather than the
Madelung pair. Then a non-integer-winding mode assigns two values to one configuration --- it is not such a
function --- and integer circulation follows by topology, with nothing added at Row~1 (conditioning on the
clock cannot manufacture a winding $\psi_0$ never had). The posit does the work a by-hand quantisation rule
used to do; it is \emph{declared}, not derived, and it is a genuine restriction --- physically multivalued
states exist (spinors, anyons) --- so the reduction is made for the simply-connected sector, and the nodal
/ multiply-connected case, where multivalued states can be physical, stays open.}
\resultpy{\text{start at }|\psi|^2\Rightarrow\text{stays }|\psi|^2\ \ (\text{equivariance})}
{DTF\_born\_relaxation\_2d.py}{H $\approx$ 0.01}
\resultpy{\text{start off-equilibrium}\Rightarrow\text{relaxes to }|\psi|^2}
{DTF\_born\_relaxation\_2d\_mp.py}{H: 0.78$\to$0.04}
\resultpy{\text{single-valued }\psi_0\ \Rightarrow\ \oint\!\nabla S\!\cdot d\ell=2\pi n\hbar\ \text{(simply-connected sector)}}
{DTF\_wallstrom\_from\_row0.py}{reduced to one posit}
\dtfresult{|\psi|^2=\text{the equilibrium the conserved location-density settles onto}}
{Born is \emph{reduced} to three premises (closing ledger), not postulated outright: it is the fixed
point of the guiding flow, reached by relaxation, with the spectrum quantised because single-valuedness
is reduced to a single Row-0 posit rather than imposed by hand. \emph{Which} outcome a single run yields --- the hard half
of measurement --- stays open. And the direction matters: this chain \emph{assumes} $\psi$ obeys the
Schr\"odinger equation and shows $\rho\to|\psi|^2$, so it \emph{reduces} Born \emph{downstream} of that
equation --- it is not a route to deriving the equation (the closing box totals what is and is not derived).}
\end{chainbox}

\chainhead{4}{Entanglement is a shared present, not a signal}
\setup{Bind quanta and you get one object with several addresses. The link is a ledger, not a channel.}
\begin{chainbox}
\known{\text{singlet: part purity } \tr\rho_A^2=0.5}{joint definiteness without individual definiteness}
\nodetext{the correlation is a Row-0 constraint (a conserved total, no location), balanced on the
\emph{proper now} --- the shared present. Separating the parts stretches nothing; measuring one reads an
entry. Nothing crosses the gap.}
\resultpy{\text{GHZ: measure one}\Rightarrow\text{other two definite \emph{or} a fresh pair; marginals unchanged}}
{DTF\_tripartite\_ledger.py}{no-signalling}
\nodetext{DTF's \emph{local} wake in ordinary space gives $\mathrm{CHSH}=2$ --- the local-beable bound
(Bell 1964~\cite{Bell1964}; Maudlin~\cite{Maudlin2011}; Norsen~\cite{Norsen2017} --- a theorem, not a misreading). To exceed it the correlation must ride the
\emph{shared present}: a joint state on one preferred slice. Relaxing that joint state to its Born
measure reproduces $\mathrm{CHSH}=2\sqrt2$ [\texttt{DTF\_bell\_from\_joint\_relaxation.py}] --- but this
is standard foliation-based relativistic Bohmian mechanics (D\"urr--Goldstein--Norsen--Zangh\`i~\cite{DGNSZ2014};
Tumulka~\cite{Tumulka2006}), and the $2\sqrt2$ is \emph{inherited} from the entangled state, not generated by the relaxation
(which only drives $\rho$ toward a given $|\psi|^2$).}
\dtfresult{\text{the shared now is the \emph{resource} for Bell violation --- and DTF gets it for free}}
{The claim is not a new derivation of $2\sqrt2$ but an \emph{ontological economy over} the foliation-based
relativistic Bohmian programme: those theories postulate the preferred foliation externally, whereas
DTF's --- the proper now --- \emph{falls out of} $\pi_u\equiv0$ (Volume~I). That is the loud, defensible
claim, and the only one made here. Whether the shared-now joint object is distinguishable from
configuration space under another name is settled, against the local reading: it is not. A general two-particle amplitude needs $O(N^2)$ functional
content and a slice carries $O(N)$, a gap no bounded internal structure closes. So the joint object
\emph{is} configuration space, and this chain claims no locality it cannot have.}
\end{chainbox}
\oneobject{the proper now --- a shared present}{The same preferred slice on which gravity solves its
field (Volume~I, Chain~7) is where an entangled system's joint state --- one object on configuration space,
not two local clocks --- is referred to a single present. One slice, two jobs; and because $c$ is universal
it stays undetectable --- no signal, no broken relativity.}

\textbf{The Schr\"odinger equation, reassembled --- and where it stops.}\enspace In polar form
$\psi=R\,e^{iS/\hbar}$ the Schr\"odinger equation splits \emph{exactly} (Madelung~\cite{Madelung1927}) into
a phase law and an amplitude law:
\[
  \underbrace{\partial_t S+\tfrac{(\nabla S)^2}{2m}+V+Q=0}_{\text{phase: Hamilton--Jacobi}},
  \qquad
  \underbrace{\partial_t R^2+\nabla\!\cdot\!\Big(R^2\,\tfrac{\nabla S}{m}\Big)=0}_{\text{amplitude: continuity}},
  \qquad Q=-\tfrac{\hbar^2}{2m}\tfrac{\nabla^2 R}{R}.
\]
\pychip{DTF\_madelung\_split.py}{the split, exhibited --- phase half native, amplitude half imported}
\par\smallskip
\textbf{The phase law is DTF-native, and reassembles from the chains above.} Its pieces are already
deposited: the accumulated clock reading $S=-mc^2\!\int u\,dt$ with its two faces $\mathbf p=\nabla S$ and
$E=-\partial_t S$ (Chain~1), and $V$ from the clock's coupling to matter (Volumes~I--II). Collected, they
are the \emph{classical} Hamilton--Jacobi equation --- missing exactly $Q$.
\par\smallskip
\textbf{The amplitude law is imported.} The continuity equation for $R^2$ is Chain~3's premise, and with it
$Q$ --- which carries \emph{all} the quantum content and is built entirely from $R$. So the quantum lives
wholly in the half DTF does not derive --- and here the framework says something definite about it.
\par\smallskip
\textbf{This is Nelson's stochastic mechanics~\cite{Nelson1966}, with the one thing Nelson left unexplained
now named.} Nelson recovers exactly this amplitude law from an underlying conservative diffusion of
coefficient $\nu=\hbar/2m$, but must \emph{posit} the diffusion; its physical origin was never supplied.
DTF supplies it --- and \emph{not} at the Row-0$\to$Row-1 interface (a separate job, below). The noise is
intrinsic to Row~1: a particle \emph{is} bound energy, and it
flexes as it oscillates inside itself at the Compton clock frequency $\omega_C=mc^2/\hbar$ --- the trembling
internal motion of a \emph{zitterbewegung}~\cite{Schrodinger1930,Hestenes1990}. That flexing is unresolvable
at the position level, and, carried into the phase $\Phi=-\omega_C\!\int u\,dt$, \emph{is} the sub-quantum
noise. So the quantum potential $Q$ is not imported machinery but the shadow of the internal flexing, and
the classical limit is the particle whose internal motion is ignorable ($Q\to0$). And because that sub-tick
flexing can \emph{never} enter a position-level deterministic formula, it is at one stroke the origin of
quantum \emph{indeterminism} --- the piece Nelson takes as brute noise.
\par\smallskip
\textbf{The scale, and the factor of $2$.} The flexing amplitude is the Compton cell $\lambda_C=\hbar/mc$,
one oscillation per tick $T_C=\hbar/mc^2$; an unbiased random walk of that step has the \emph{standard}
diffusion coefficient
\[ \nu\;=\;\frac{\lambda_C^{2}}{2\,T_C}\;=\;\frac{\hbar}{2m}, \]
which is Nelson's $\nu$ exactly --- though the factor of $2$ carries the one-dimensional random-walk
convention $\langle x^2\rangle=2\nu t$, and is not itself fixed by the flexing dynamics. So $\hbar$ is
\emph{reparented}: not posited afresh as ``the quantum of action'' but read as the scale of the particle's
own internal flexing. Writing the step as $\alpha\lambda_C$ ($\nu=\alpha^2\hbar/2m$) and requiring the
osmotic potential it builds to \emph{be} the Madelung potential $Q$,
\[ -2m\nu^2\,\frac{\nabla^2 R}{R}\ =\ -\frac{\hbar^2}{2m}\,\frac{\nabla^2 R}{R}, \]
returns $\alpha=1$: a \emph{consistency check} that the two roles of $\hbar$ agree, not an independent
derivation of the coefficient. What remains open --- and is listed as owed below --- is whether the flexing
dynamics \emph{forces} $\hbar/2m$ rather than merely admitting it.
\pychip{DTF\_nelson\_from\_flexing.py}{$\nu=\hbar/2m$ (1D convention); osmotic-match returns $\alpha{=}1$ --- consistency, not derivation}
\par\smallskip
\textbf{Conservative, not dissipative --- and that is the whole of it.} Nelson's construction yields
Schr\"odinger rather than a heat equation only if the diffusion is \emph{conservative} (time-symmetric),
his most-criticised assumption. Two DTF facts point the same way, and a free-particle check confirms the
consequence: the internal flexing is a \emph{bound} oscillation --- rest mass neither radiates nor decays
--- and a bound, finite motion cannot dissipate (dissipation needs a thermal continuum), so it is
reversible; and Row~0 is timeless, with no substrate arrow. For the free particle the time-symmetric
(Nelson) acceleration vanishes \emph{exactly} when the osmotic coefficient equals $\hbar/2m$, reproducing
quantum spreading, whereas the forward-only (arrow-ful) rule gives the heat equation
[\texttt{DTF\_nelson\_from\_flexing.py}]. Whether the flexing dynamics realises this \emph{exactly} is the
sector's one open item --- stated as its single posit in the box below.
\par\smallskip
\textbf{The interface, separately, carries conservation and single-valuedness.} Distinct from the Row-1
flexing, the Row-0$\to$Row-1 conditioning is what supplies the continuity law (Chain~3) and preserves
single-valuedness of $\psi$ (Chain~3's Wallstrom posit --- it manufactures no winding). Noise on the
amplitude; conservation and topology at the interface: two different jobs, not to be run together.
\smallskip
\noindent\textbf{At true size:} DTF does not \emph{derive} the Schr\"odinger equation. It supplies the
phase half structurally and \emph{reduces} the amplitude half to a single physical postulate --- the
conservative internal flexing of Row-1 bound energy --- with the coefficient now native ($\hbar/2m$, the
random-walk normalisation) and the flexing's time-symmetry now shown to follow from bound-energy and a
timeless Row~0 (harmonic model; only full nonlinear rigor remains --- see the posit box). And \emph{reducing Born
is not deriving Schr\"odinger}: Chain~3 assumes $\psi$ obeys the equation and shows $\rho\to|\psi|^2$ ---
downstream of it. What is new is an \emph{identification}: Nelson's unexplained noise, and quantum
indeterminism with it, is the internal flexing of bound energy --- the always-on, $\hbar$-scale form of the
same clock indefiniteness whose growth to an $O(1)$ crossing (there gravitationally driven) is measurement
and the black-hole core. The trilogy marks the wall exactly here.

\textbf{The one quantum posit --- weighted with the singularity's.}\enspace Volume~II's black hole rests on
a single ontological posit: Row-1 time \emph{de-emerges} where the clock loses definiteness
($\delta u/u\sim1$), and the finite core, the Planck remnant, and the absent echoes follow. The quantum
sector rests on one of the same weight. Given what the trilogy has already fixed --- the density
$\rho=|\psi|^2$, the phase $S$ as the clock reading, and the flexing scale $\nu=\hbar/2m$ --- the
variational principle for the amplitude law (Madelung; Guerra--Morato~\cite{GuerraMorato1983}) has
\emph{one} free choice left: the sign of its osmotic term,
\[ \varepsilon=+1\ \Rightarrow\ \text{real-time, unitary: Schr\"odinger};\qquad
   \varepsilon=-1\ \Rightarrow\ \text{imaginary-time, dissipative: a heat equation}. \]
\pychip{DTF\_qm\_posit\_variational.py}{$\varepsilon{=}{+}1$ Euler--Lagrange $\Leftrightarrow$ Schr\"odinger (verified); $\varepsilon{=}{-}1\Rightarrow$ heat}
That sign \emph{is} the time-symmetry of the internal flexing --- a flexing with no preferred time
direction takes the real-time branch. So the whole quantum sector reduces to one statement:
\begin{center}\itshape The internal flexing of bound energy is time-symmetric.\end{center}
\noindent And $\varepsilon=+1$ is not an arbitrary choice: \emph{within a harmonic model of the flexing} it
follows from what DTF holds. Model the flexing as its \emph{own} (internal, not environmental)
Caldeira--Leggett~\cite{CaldeiraLeggett1983} sector: a \emph{bound}, discrete oscillator has a vanishing
friction kernel ($\tilde\gamma(s\!\to\!0)=0$, reactive only) and returns the coarse energy in \emph{full}
(reversible), so the induced diffusion is time-symmetric; and Row~0, being timeless, imports no arrow.
\emph{One caveat, stated plainly:} the physical flexing is a \emph{zitterbewegung}, which in Dirac theory
is interference across the positive/negative-energy \emph{continuum}, not a few discrete modes --- so
whether it retains that time-symmetric character is part of the relativistic completion listed as open
below, not settled here. Within the harmonic model, a vanishing friction kernel \emph{is} a
time-symmetric influence action, which \emph{is} $\varepsilon=+1$
[\texttt{DTF\_flexing\_time\_symmetry.py}]. So the quantum sector's character is a \emph{consequence} of the
bound-energy ontology and a timeless substrate --- shown here for the harmonic flexing model. It sits at the
other end of the same axis as the singularity: the singularity is the \emph{ceiling} of clock indefiniteness
($\delta u/u\to1$, time de-emerges), the flexing its always-on \emph{floor} ($\hbar$-scale, definite but
never perfectly). One knob, two ends. \textbf{The honest remainder} is the relativistic one: whether the
physical flexing --- a continuum zitterbewegung, not the harmonic model's discrete mode --- retains the
time-symmetric character; that is the relativistic-completion question below, not settled here.

\begin{tcolorbox}[keybox]
\textbf{At relativistic speed: an open program, and where DTF stands.}\enspace The amplitude law above is
non-relativistic (Schr\"odinger), while the phase is relativistic ($\Phi=-\omega_C\!\int d\tau$, with
$\mathbf p=\gamma m\mathbf v$). Closing that gap is an open program --- but an unusually favourable one here,
for two structural reasons. \emph{First}, relativistic stochastic mechanics classically fails for want of a
preferred time in which the diffusion runs; DTF already has one --- the proper now ($\pi_u\equiv0$), the
same slice that does gravity and entanglement --- so the natural covariant statement is that the diffusion
runs in \emph{proper time}, $\nu=\hbar/2m$ per $d\tau$, exactly as the phase does. \emph{Second}, the noise
is a \emph{zitterbewegung}, a Dirac-equation phenomenon; so the relativistic amplitude sector should
complete not toward a patched Schr\"odinger but toward \emph{Dirac} (Hestenes~\cite{Hestenes1990}), with
Schr\"odinger as its $\gamma\to1$ face. \emph{The honest ceiling}: the flexing amplitude \emph{is} the
Compton length $\lambda_C$ --- precisely the scale at which single-particle quantum mechanics gives way to
pair production and quantum field theory. So the single-particle construction is intrinsically bounded below
the Compton scale, and the full completion is ultimately field-theoretic, not merely Dirac. None of this is
claimed here; it is the sharpest forward direction, and DTF holds the two ingredients the direction needs
--- a preferred time, and a relativistic noise.
\end{tcolorbox}

\textbf{What is earned, borrowed, declared, and owed.}
\begin{itemize}[leftmargin=1.4em,itemsep=2pt,topsep=3pt]
  \item \textbf{Earned.} The phase as a clock reading; the one phase $\Phi$ as the shared root of inertia
    (the rate), gravity (the well it sources), and guidance (its space-face gradient); the amplitude as
    the same clock on the space axis; entanglement as
    joint-without-individual definiteness on the proper now; \textbf{the Born rule}, as the relaxation
    equilibrium of the conserved location-density (verified), with quantised circulation
    \emph{reduced} to a single Row-0 posit (a single-valued amplitude assignment) rather than imposed by
    hand --- for the simply-connected sector.
  \item \textbf{Borrowed.} The Tsirelson bound $2\sqrt2$~\cite{Cirelson1980}; the sub-quantum $H$-theorem
    (Valentini~\cite{Valentini1991}) and the diffusion construction of the amplitude sector
    (Nelson~\cite{Nelson1966}).
  \item \textbf{Retired --- not borrowed.} The measurement-axiom scaffolding
    (Gleason~\cite{Gleason1957}, Busch~\cite{Busch2003}) is not inherited here but \emph{dropped}: those
    theorems justify $|\psi|^2$ as the unique measure over measurement outcomes, and once $|\psi|^2$ is
    instead the equilibrium of a conserved location density, that is a question this account does not
    ask. Stated as its own commitment rather than a clause, because it is the stronger claim: not that
    the scaffolding is used differently, but that it is unnecessary.
  \item \textbf{Declared.} The substrate is unresolved to any observer built from it (typicality);
    correlations are overlaps (exact for the qubit); and the posit that carries the Wallstrom reduction ---
    that Row~0's beable is a \emph{single-valued} amplitude assignment on configuration space. This is what
    replaces a by-hand quantisation rule; it is a genuine restriction (spinors and anyons are physically
    multivalued), so it is declared for the simply-connected sector, not derived in general.
  \item \textbf{Owed.} Coarse-graining (typicality); and \emph{which} single outcome a run yields
    (selection), which every no-collapse ontology owes. \emph{The amplitude sector's dynamics:} the phase
    half of the Schr\"odinger equation is DTF-native (the reassembly above), but the amplitude half is
    imported. For its origin DTF offers a \emph{candidate}, not a derivation --- the clock's cell-to-cell
    resynchronisation, whose fluctuation is a natural source for the diffusion the amplitude law needs ---
    at a scale $\hbar$ posited as every formulation of quantum mechanics posits it; what is open is whether
    that dynamics \emph{forces} the coefficient $\hbar/2m$ (the factor of $2$). DTF therefore \emph{reduces}
    the Schr\"odinger equation to one posited constant plus one open coefficient; it does not derive it.
    \emph{On single-valuedness, two open items, honestly sized:} (i)~that the Row-0$\to$Row-1 emergence
    through the $u$-clock is a conditioning-type map (which would preserve single-valuedness), still to be
    shown for the specific $u$-clock emergence; and (ii)~the nodal / multiply-connected case, genuinely
    open, because there the single-valuedness posit is a real restriction (spinors, anyons). These are
    smaller than ``Wallstrom is unaddressed'' --- the obstruction is reduced to one declared posit --- but
    larger than a single technical residual.
  \item \textbf{Settled, negatively.} The Bell question of Chain~4 --- whether the shared-now joint object
    is distinguishable from a \emph{fundamental} configuration space --- is not owed: it has
    been answered, and it is not. The guiding wave is on configuration space, and this volume claims no
    locality beyond the single quantum.
\end{itemize}

\noindent On gravity-mediated entanglement (BMV): \emph{if} the clock-rate constraint is quantised ---
operator-valued, as QED's non-radiative Coulomb sector is --- the guiding-not-driving reading predicts a
\emph{positive} outcome. This is stated conditionally, and honestly: the zero-DOF and polarisation
results (Volume~I--II) treat $u$ as a \emph{classical} constraint, whereas BMV needs it operator-valued
--- a quantisation this trilogy does not construct. So DTF's BMV stance is ``positive, conditional on
quantising the constraint,'' not a delivered prediction; the framework's delivered falsifiable stakes are
in Volume~II (GW polarisations, black-hole cores).

\medskip
\noindent The reading offered here is not a separate theory of measurement but the \emph{organizing
principle} of Volume~I holding in the quantum sector: ``becoming definite'' is the Row-0$\to$Row-1 step ---
the \emph{same} clock-definiteness surface whose opposite crossing is the black-hole core of Volume~II,
where a clock instead de-emerges. Measurement and the singularity are one transition, run in the two
directions; that is the sense in which the three volumes are one principle, not three results. A fourth ---
deep cosmology (the proper now as cosmic time) --- is in preparation.

\medskip
\noindent{\small\raggedright Scripts (public, SymPy/NumPy): \texttt{DTF\_two\_clocks\_resolved},
\texttt{DTF\_guidance\_is\_spaceface}, \texttt{DTF\_phase\_amplitude\_oneconstant},
\texttt{DTF\_doubleslit\_compton\_ticks}, \texttt{DTF\_born\_relaxation\_2d},
\texttt{DTF\_born\_relaxation\_2d\_mp}, \texttt{DTF\_wallstrom\_from\_row0},
\texttt{DTF\_sharednow\_counting}, \texttt{DTF\_tripartite\_ledger},
\texttt{DTF\_bell\_from\_joint\_relaxation}, \texttt{DTF\_clockjitter\_scale},
\texttt{DTF\_nelson\_from\_flexing}, \texttt{DTF\_qm\_posit\_variational},
\texttt{DTF\_flexing\_time\_symmetry}.\par}

\begin{tcolorbox}[keybox]
\textbf{Provenance of the checks (ontology-first).}\enspace This volume derives little that is new ---
that is the point. Its contribution is the \emph{reading}, not new physics, and the scripts are labelled
accordingly. \textbf{Borrowed} (standard constructions, credited): the relaxation to $|\psi|^2$
(Valentini's $H$-theorem on Bohm's guidance) and its joint-configuration extension to $2\sqrt2$; the
stochastic-quantisation and Nelson checks. \textbf{Consistency checks} (a standard fact or definition
confirmed, \emph{not} a native derivation): the two-clock $c$-number/state split, the one-constant
Planck/Compton identity, the double-slit Compton-tick $=$ de~Broglie identity, the singlet purity and
Horodecki CHSH~\cite{Horodecki1995}, the GHZ no-signalling. Two \textbf{DTF-native structural} checks
earn their place: $u\equiv1\Rightarrow\nabla S\neq0$ (\texttt{DTF\_guidance\_is\_spaceface}), which
\emph{refutes} the reading of pilot-wave guidance as a wake in the clock rate, fixing guidance as the
space face of the interval; and the Wallstrom reduction
(\texttt{DTF\_wallstrom\_from\_row0}) --- that single-valuedness reduces to one declared Row-0 posit (a
single-valued amplitude assignment) rather than a by-hand quantisation rule, for the simply-connected
sector. (The circulation Wallstrom asks about is the \emph{spatial} configuration phase on a fixed slice,
not the unbounded temporal one, so no clock-angle argument bears on it; that framing is retired.) The bulk of the native computational content still lives in Volume~II; here the work is
ontological. None of it rests on circular reasoning.
\end{tcolorbox}

\appendix
\section*{Appendix A --- The DTF Core in brief}
\addcontentsline{toc}{section}{Appendix A}
\noindent For readers coming to this volume first. The Distorted Time Framework takes the local
\emph{rate} of time as the one physical primitive, and builds outward in three layers (Volume~I).

\begin{itemize}[leftmargin=1.4em,itemsep=3pt,topsep=3pt]
  \item \textbf{Row 0 --- the timeless substrate.} The deepest layer has no clock: the Wheeler--DeWitt
    constraint $\hat H\Psi=0$ holds with no external time. Correlations are simultaneous, not ordered.
  \item \textbf{Row 1 --- the clock.} A definite local clock-rate field $u(\mathbf x)>0$ is conditioned
    on Row~0; its reading $\tau=\int u\,dt$ accumulates monotonically --- time, and its arrow. A
    quantum's phase is this reading, $\Phi=-\omega_C\tau$.
  \item \textbf{Row 2 --- space.} An extent is a period radiated at $c$ ($\lambda_C=c\,T_C$): space is
    time's radiation, compiled by an orientation (frame) field. Because $c$ is one universal exchange
    rate, Lorentz invariance is its signature.
  \item \textbf{One object, four faces.} Extremising the interval gives one phase read four ways: inertia
    ($\mathbf p=\nabla S$), free fall (geodesic), the pilot wave ($\mathbf v=\nabla S/m$), and the path
    integral. Inertia (the rate $\omega_C$) and the pilot wave (the space-face gradient $\nabla S$) are
    tied by $\mathbf p=\nabla S=m\mathbf v$, both $\hbar$-set and linear in $m$; the $G$-coupled well the
    rate digs is gravity, a separate object.
  \item \textbf{The proper now.} The clock rate does not propagate ($\pi_u\equiv0$), so it is fixed
    across space on one preferred slice --- a real but empirically inert present, shared by gravity and
    by entanglement.
  \item \textbf{Definiteness.} ``Becoming definite'' is the Row-0$\to$1 step; a single
    clock-definiteness threshold governs measurement (this volume) and the black-hole core (Volume~II).
\end{itemize}

\section*{Code and data availability}
\addcontentsline{toc}{section}{Code and data availability}
{\small The SymPy/NumPy verification scripts named throughout (the \texttt{DTF\_born\_relaxation} and
\texttt{DTF\_wallstrom\_from\_row0} scripts, \texttt{DTF\_nelson\_from\_flexing},
\texttt{DTF\_qm\_posit\_variational}, \texttt{DTF\_flexing\_time\_symmetry}, and the entanglement/counting
scripts) are publicly available (see the repository link on the title page) at the release tagged
\texttt{v1.0} (Version~8 of this trilogy); see the repository's README for a script-to-claim index and run
instructions.\par}

\section*{Author's Note: the role of AI}
\addcontentsline{toc}{section}{Author's Note: the role of AI}
{\small
The Distorted Time Framework---its central hypothesis (time's local \emph{rate}, the ADM lapse,
as the one physical primitive), its ontology-first strategy, and the physical
interpretations---was conceived and developed by the author over several years of independent
research outside academic institutions. The mathematical formalization was carried out in
collaboration with Claude (Anthropic): the Born-rule relaxation and Wallstrom-reduction arguments,
the entanglement, Bell, and functional-content counting results, and the symbolic/numerical
verification scripts (SymPy/NumPy) referenced throughout. Separately, other large language models
(ChatGPT, Grok, Gemini) were used as independent referees, providing adversarial review of the
manuscript; their critiques shaped the revisions but not the derivations.

In each case the author specified the physical target, the constraints the result had to satisfy,
and the direction of the argument, and reviewed and approved each result; the AI executed and
checked the formal derivations. Without that assistance the author could not have produced the
mathematical formalization at this level of detail.

This is disclosed because intellectual honesty requires it. It does not change the epistemic status
of the results---a derivation is correct or it is not, independent of how it was produced---but it
does mean they warrant careful independent scrutiny rather than assumed authority. The author
welcomes that scrutiny; the explicit claim-type labels and the reproducible verification scripts are
provided precisely to enable it.
\par}

\begingroup\small
\begin{thebibliography}{99}
\addcontentsline{toc}{section}{References}
\bibitem{Born1926} M.~Born, ``Zur Quantenmechanik der Sto\ss vorg\"ange,''
  \emph{Z.\ Phys.} \textbf{37}, 863--867 (1926).
\bibitem{Nelson1966} E.~Nelson, ``Derivation of the Schr\"odinger Equation from Newtonian Mechanics,''
  \emph{Phys.\ Rev.} \textbf{150}, 1079--1085 (1966).
\bibitem{Wallstrom1994} T.~C.~Wallstrom, ``Inequivalence between the Schr\"odinger equation and the
  Madelung hydrodynamic equations,'' \emph{Phys.\ Rev.\ A} \textbf{49}, 1613--1617 (1994).
\bibitem{Bohm1952} D.~Bohm, ``A Suggested Interpretation of the Quantum Theory in Terms of `Hidden'
  Variables. I \& II,'' \emph{Phys.\ Rev.} \textbf{85}, 166--179, 180--193 (1952).
\bibitem{Bell1964} J.~S.~Bell, ``On the Einstein Podolsky Rosen Paradox,''
  \emph{Physics} \textbf{1}, 195--200 (1964).
\bibitem{Maudlin2011} T.~Maudlin, \emph{Quantum Non-Locality and Relativity}, 3rd ed.
  (Wiley-Blackwell, 2011).
\bibitem{Norsen2017} T.~Norsen, \emph{Foundations of Quantum Mechanics} (Springer, 2017).
\bibitem{DGNSZ2014} D.~D\"urr, S.~Goldstein, T.~Norsen, W.~Struyve, and N.~Zangh\`i, ``Can Bohmian
  mechanics be made relativistic?'' \emph{Proc.\ R.\ Soc.\ A} \textbf{470}, 20130699 (2014).
\bibitem{Tumulka2006} R.~Tumulka, ``A Relativistic Version of the Ghirardi--Rimini--Weber Model,''
  \emph{J.\ Stat.\ Phys.} \textbf{125}, 821--840 (2006).
\bibitem{Madelung1927} E.~Madelung, ``Quantentheorie in hydrodynamischer Form,''
  \emph{Z.\ Phys.} \textbf{40}, 322--326 (1927).
\bibitem{Schrodinger1930} E.~Schr\"odinger, ``\"Uber die kr\"aftefreie Bewegung in der relativistischen
  Quantenmechanik'' (Zitterbewegung), \emph{Sitzungsber.\ Preuss.\ Akad.\ Wiss.\ Phys.-Math.\ Kl.}
  \textbf{24}, 418--428 (1930).
\bibitem{Hestenes1990} D.~Hestenes, ``The Zitterbewegung Interpretation of Quantum Mechanics,''
  \emph{Found.\ Phys.} \textbf{20}, 1213--1232 (1990).
\bibitem{GuerraMorato1983} F.~Guerra and L.~M.~Morato, ``Quantization of dynamical systems and stochastic
  control theory,'' \emph{Phys.\ Rev.\ D} \textbf{27}, 1774--1786 (1983).
\bibitem{CaldeiraLeggett1983} A.~O.~Caldeira and A.~J.~Leggett, ``Path integral approach to quantum
  Brownian motion,'' \emph{Physica A} \textbf{121}, 587--616 (1983).
\bibitem{Cirelson1980} B.~S.~Cirel'son (Tsirelson), ``Quantum generalizations of Bell's inequality,''
  \emph{Lett.\ Math.\ Phys.} \textbf{4}, 93--100 (1980).
\bibitem{Valentini1991} A.~Valentini, ``Signal-locality, uncertainty, and the subquantum
  $H$-theorem. I,'' \emph{Phys.\ Lett.\ A} \textbf{156}, 5--11 (1991).
\bibitem{Gleason1957} A.~M.~Gleason, ``Measures on the Closed Subspaces of a Hilbert Space,''
  \emph{J.\ Math.\ Mech.} \textbf{6}, 885--893 (1957).
\bibitem{Busch2003} P.~Busch, ``Quantum States and Generalized Observables: A Simple Proof of Gleason's
  Theorem,'' \emph{Phys.\ Rev.\ Lett.} \textbf{91}, 120403 (2003).
\bibitem{Horodecki1995} R.~Horodecki, P.~Horodecki, and M.~Horodecki, ``Violating Bell inequality by
  mixed spin-$\tfrac12$ states,'' \emph{Phys.\ Lett.\ A} \textbf{200}, 340--344 (1995).
\end{thebibliography}
\endgroup

\end{document}
